% Part: model-theory
% Chapter: interpolation
% Section: interpolation-proof

\documentclass[../../../include/open-logic-section]{subfiles}

\begin{document}

\olfileid{mod}{int}{prf}

\section{ક્રેગનું અંતર્વેશન પ્રમેય}

\begin{thm}[ક્રેગનું અંતર્વેશન પ્રમેય]
\ollabel{thm:interpol}
જો $\Entails !A \lif !B$, તો એવું !!a{sentence}~$!C$ છે કે
$\Entails !A \lif !C$ અને $\Entails !C \lif !B$, અને $!C$માં આવતો
દરેક !!{constant}, !!{function} અને !!{predicate} ($\eq$ સિવાયનો)
$!A$ અને~$!B$ બંનેમાં આવે છે. !!{sentence}~$!C$ને $!A$ અને~$!B$નો
\emph{અંતર્વેશક} કહે છે.
\end{thm}

\begin{proof}
ધારો કે $\Lang{L}_1$, $!A$ની ભાષા છે અને $\Lang{L}_2$, $!B$ની ભાષા
છે. $\Lang{L}_0 = \Lang{L}_1 \cap \Lang{L}_2$ રાખો. દરેક $i \in
\{0, 1, 2 \}$ માટે, $\Lang{L}_i$માં અનંત નવા !!{constant}s
$\Obj c_0, \Obj c_1, \Obj c_2, \dots$ ઉમેરીને $\Lang{L}'_i$ મેળવો.

જો $!A$ અસંતોષ્ય હોય, તો $\lexists[x][\eq/[x][x]]$ અંતર્વેશક છે.
જો $\lnot !B$ અસંતોષ્ય હોય (અને તેથી $!B$ પ્રામાણ્ય હોય), તો
$\lexists[x][\eq[x][x]]$ અંતર્વેશક છે. તેથી આપણે $!A$ અને
$\lnot !B$ બંને સંતોષ્ય છે એમ પણ ધારી શકીએ છીએ.

અંતર્વેશન પ્રમેયનું પ્રતિપક્ષી સ્વરૂપ સાબિત કરવા માટે ધારો કે $!A$
અને~$!B$ માટે કોઈ અંતર્વેશક નથી. બીજા શબ્દોમાં, ધારો કે $\{!A\}$ અને
$\{\lnot !B\}$, $\Lang{L}_0$માં અપૃથક્કરણીય છે.

આપણો ઉદ્દેશ જોડ $(\{ !A \}, \{\lnot!B\})$ને કોઈ મહત્તમ
અપૃથક્કરણીય જોડ $(\Gamma^*, \Delta^*)$ સુધી વિસ્તારવાનો છે.
$!A_0$, $!A_1$, $!A_2$,~\dots એ $\Lang{L}'_1$નાં !!{sentence}sની
ગણના અને $!B_0$, $!B_1$, $!B_2$,~\dots એ $\Lang{L}'_2$નાં
!!{sentence}sની ગણના હોય. નીચે પ્રમાણે $n\ge0$ માટે !!{sentence}sના
ગણોની બે વધતી શ્રેણીઓ $(\Gamma_n,\Delta_n)$ વ્યાખ્યાયિત કરીએ. $\Gamma_0
= \{!A\}$ અને $\Delta_0=\{\lnot !B\}$ રાખો. $(\Gamma_n,\Delta_n)$
પહેલેથી વ્યાખ્યાયિત છે એમ ધારીને, $\Gamma_{n+1}$ અને
$\Delta_{n+1}$ને આ પ્રમાણે વ્યાખ્યાયિત કરો:
\sourcecorrection{OLINT-003}{સ્થિર અંગ્રેજી મૂળની
  \texttt{interpolation-proof.tex}, પંક્તિઓ~41--43માં $!A_n$ અને
  $!B_n$ને અનુક્રમે $\Lang{L}_1$ અને $\Lang{L}_2$નાં વાક્યોની ગણના
  કહેવાય છે. પરંતુ પછીની પંક્તિઓ~107--108માં રચાયેલા ગણો
  $\Lang{L}'_1$ અને $\Lang{L}'_2$માં પૂર્ણ સુસંગત હોવાનો દાવો થાય છે,
  જેના માટે નવા અચળો ધરાવતાં વાક્યોની પણ ગણના જરૂરી છે. અહીં બંને
  વિસ્તૃત ભાષાઓ પર ગણના પુનઃસ્થાપિત કરી છે.}
\begin{enumerate}
\item જો $\Gamma_n \cup \{!A_n \}$ અને~$\Delta_n$, $\Lang{L}'_0$માં
  અપૃથક્કરણીય હોય, તો $!A_n$ને $\Gamma_{n+1}$માં મૂકો. વધુમાં, જો
  $!A_n$ કોઈ અસ્તિત્વવાચક !!{formula}~$\lexists[x][!S]$ હોય, તો
  $\Gamma_n$, $\Delta_n$, $!A_n$ કે~$!B_n$માં ન આવતો નવો
  !!{constant}~$c$ પસંદ કરીને $\Subst{!S}{c}{x}$ને
  $\Gamma_{n+1}$માં મૂકો.
\item જો $\Gamma_{n+1}$ અને $\Delta_n \cup \{!B_n \}$,
  $\Lang{L}'_0$માં અપૃથક્કરણીય હોય, તો $!B_n$ને $\Delta_{n+1}$માં
  મૂકો. વધુમાં, જો $!B_n$ કોઈ અસ્તિત્વવાચક
  !!{formula}~$\lexists[x][!S]$ હોય, તો $\Gamma_{n+1}$, $\Delta_n$,
  $!A_n$ કે~$!B_n$માં ન આવતો નવો !!{constant}~$c$ પસંદ કરીને
  $\Subst{!S}{c}{x}$ને $\Delta_{n+1}$માં મૂકો.
\end{enumerate}
છેવટે, આ પ્રમાણે વ્યાખ્યા કરો:
\begin{align*}
  \Gamma^* & = \bigcup_{n\ge 0} \Gamma_n, &
  \Delta^* & = \bigcup_{n\ge 0} \Delta_n.
\end{align*}
હવે $n$ પરના સમકાલીન આગમનથી આપણે નીચેના દાવા સાબિત કરી શકીએ છીએ:
\begin{enumerate}
\item\ollabel{part-a} $\Gamma_n$ અને~$\Delta_n$, $\Lang{L}'_0$માં
  અપૃથક્કરણીય છે;
\item\ollabel{part-b} $\Gamma_{n+1}$ અને~$\Delta_n$,
  $\Lang{L}'_0$માં અપૃથક્કરણીય છે.
\end{enumerate}
\olref{part-a}નો આધાર \olref[sep]{lem:sep1} આપે છે. ભાગ
\olref{part-b} માટે આપણે ત્રણ કિસ્સા પાડવા પડે છે:
\begin{enumerate}
\item જો $\Gamma_0 \cup \{!A_0 \}$ અને~$\Delta_0$ પૃથક્કરણીય હોય,
  તો $\Gamma_1=\Gamma_0$ અને \olref{part-b}, \olref{part-a} જ છે;
\item જો $\Gamma_1 = \Gamma_0 \cup\{!A_0\}$, તો રચના મુજબ
  $\Gamma_1$ અને~$\Delta_0$ અપૃથક્કરણીય છે.
\item હવે માત્ર $!A_0$ અસ્તિત્વવાચક હોય તે કિસ્સો બાકી રહે છે, જેમાં
  $\Gamma_1 = \Gamma_0 \cup \{ \lexists[x][!S], \Subst{!S}{c}{x}
  \}$. રચના મુજબ $\Gamma_0 \cup \{ \lexists[x][!S]\}$ અને
  $\Delta_0$ અપૃથક્કરણીય છે, તેથી \olref[sep]{lem:sep2} મુજબ
  $\Gamma_0 \cup \{ \lexists[x][!S], \Subst{!S}{c}{x} \}$ અને
  $\Delta_0$ પણ અપૃથક્કરણીય છે.
\end{enumerate}
આથી ઉપરના \olref{part-a} અને \olref{part-b} માટે આગમનનો આધાર પૂરો
થાય છે. હવે આગમન-પગલું લઈએ. \olref{part-a} માટે, જો $!B_n$ને
$\Delta_{n+1}$માં મૂકવામાં આવે, તો રચના મુજબ $\Gamma_{n+1}$ અને
$\Delta_n\cup\{!B_n\}$ અપૃથક્કરણીય છે; અને જો $!B_n$ અસ્તિત્વવાચક
હોય, તો સાક્ષી ઉમેર્યા પછીનો $\Delta_{n+1}$ પણ
\olref[sep]{lem:sep2} મુજબ અપૃથક્કરણીય રહે છે. જો
$\Delta_{n+1}=\Delta_n$ (કારણ કે $\Gamma_{n+1}$ અને
$\Delta_n\cup\{!B_n\}$ પૃથક્કરણીય છે), તો \olref{part-b} માટેની
આગમન-પરિકલ્પના વાપરીએ. \olref{part-b}ના આગમન-પગલા માટે, જો
$!A_{n+1}$ને $\Gamma_{n+2}$માં મૂકવામાં આવે, તો રચના મુજબ
$\Gamma_{n+1}\cup\{!A_{n+1}\}$ અને~$\Delta_{n+1}$ અપૃથક્કરણીય છે;
અને જો $!A_{n+1}$ અસ્તિત્વવાચક હોય, તો સાક્ષી ઉમેર્યા પછીનો
$\Gamma_{n+2}$ પણ \olref[sep]{lem:sep2} મુજબ અપૃથક્કરણીય રહે છે. જો
$\Gamma_{n+2}=\Gamma_{n+1}$, તો હમણાં જ સાબિત કરેલો
\olref{part-a}નો આગમન-કિસ્સો વાપરીએ. આથી \olref{part-a} અને
\olref{part-b} પરનું આગમન પૂરું થાય છે.
\sourcecorrection{OLINT-004}{સ્થિર અંગ્રેજી મૂળની
  \texttt{interpolation-proof.tex}, પંક્તિઓ~88--96માં અસ્તિત્વવાચક
  $!B_n$ કે $!A_{n+1}$ સ્વીકારાય ત્યારે સાક્ષી-સૂત્ર પણ ઉમેરવાની
  અગાઉની સૂચના હોવા છતાં અનુક્રમે
  $\Delta_{n+1}=\Delta_n\cup\{!B_n\}$ અને
  $\Gamma_{n+2}=\Gamma_{n+1}\cup\{!A_{n+1}\}$ લખાયું છે. અહીં
  સ્વીકારેલો વિધાન અને પછી ઉમેરાતો સાક્ષી અલગ કરીને, લેમા
  \olref[sep]{lem:sep2}નો ઉપયોગ બંને પૂર્ણ ગણો પર સ્પષ્ટ કર્યો છે.}

આમાંથી ફલિત થાય છે કે $\Gamma^*$ અને~$\Delta^*$ અપૃથક્કરણીય છે;
નહિતર સઘનતાથી એવું $n\ge0$ અને એવું $\Lang{L}'_0$-વિધાન મળે જે
$\Gamma_n$ અને~$\Delta_n$ને પૃથક કરે, જે \olref{part-a}ની વિરુદ્ધ
છે. ખાસ કરીને, $\Gamma^*$ અને~$\Delta^*$ સુસંગત છે: કારણ કે પહેલો
અથવા બીજો અસુસંગત હોય, તો અનુક્રમે $\lexists[x][\eq/[x][x]]$ અથવા
$\lforall[x][\eq[x][x]]$ તેમને પૃથક કરે.
\sourcecorrection{OLINT-005}{સ્થિર અંગ્રેજી મૂળની
  \texttt{interpolation-proof.tex}, પંક્તિઓ~100--102માં કોઈ
  $n\ge0$ પોતે $\Gamma_n$ અને $\Delta_n$ને પૃથક કરે છે તેમ લખાયું
  છે. પૃથક્કર્તા વિધાન હોય છે; સઘનતા માત્ર એવું એક જ~$n$ આપે છે જેમાં
  તેના માટે જરૂરી બંને સાન્ત ઉપગણો આવી જાય. અહીં એ વિધાન અને~$n$ની
  જુદી ભૂમિકાઓ સ્પષ્ટ કરી છે.}

હવે આપણે બતાવીએ કે $\Gamma^*$, $\Lang{L}'_1$માં પૂર્ણ સુસંગત છે અને
એ જ રીતે $\Delta^*$, $\Lang{L}'_2$માં પૂર્ણ સુસંગત છે. પહેલા માટે,
ધારો કે કોઈ $n\ge0$ માટે $!A_n\notin\Gamma^*$ અને $\lnot !A_n
\notin\Gamma^*$. જો $!A_n\notin\Gamma^*$, તો $\Gamma_n\cup\{!A_n\}$
$\Delta_n$થી પૃથક્કરણીય છે, તેથી એવું $!C\in\Lang{L}'_0$ છે કે આ
બંને વાતો સાચી છે:
\begin{align*}
  \Gamma^* & \Entails !A_n \lif !C, &
  \Delta^* & \Entails \lnot !C.
\end{align*}
એ જ રીતે, જો $\lnot !A_n\notin\Gamma^*$, તો એવું $!C'\in
\Lang{L}'_0$ છે કે આ બંને વાતો સાચી છે:
\begin{align*}
  \Gamma^* & \Entails \lnot !A_n \lif !C', &
  \Delta^* & \Entails \lnot !C'.
\end{align*}
વિધાનાત્મક તર્કશાસ્ત્રથી $\Gamma^*\Entails !C\lor !C'$ અને
$\Delta^*\Entails\lnot(!C\lor !C')$, તેથી $!C\lor !C'$,
$\Gamma^*$ અને~$\Delta^*$ને પૃથક કરે છે. સમાન દલીલથી $\Delta^*$ પૂર્ણ
છે.

છેવટે, આપણે બતાવીએ કે $\Gamma^*\cap\Delta^*$, $\Lang{L}'_0$માં પૂર્ણ
સુસંગત છે. તે સુસંગત ગણોનો છેદ હોવાથી દેખીતી રીતે સુસંગત છે. પૂર્ણતા
બતાવવા $!S\in\Lang{L}'_0$ લો. હવે $\Gamma^*$, $\Lang{L'_1}
\supseteq\Lang{L'_0}$માં પૂર્ણ છે અને એ જ રીતે $\Delta^*$,
$\Lang{L'_2}\supseteq\Lang{L'_0}$માં પૂર્ણ છે. તેથી કાં તો
$!S\in\Gamma^*$ અથવા $\lnot !S\in\Gamma^*$, અને કાં તો
$!S\in\Delta^*$ અથવા $\lnot !S\in\Delta^*$. જો $!S\in\Gamma^*$ અને
$\lnot !S\in\Delta^*$, તો $!S$, $\Gamma^*$ અને~$\Delta^*$ને પૃથક
કરે; અને જો $\lnot !S\in\Gamma^*$ અને $!S\in\Delta^*$, તો
$\Gamma^*$ અને~$\Delta^*$ને $\lnot !S$ પૃથક કરે. આથી, કાં તો
$!S\in\Gamma^*\cap\Delta^*$
અથવા $\lnot !S\in\Gamma^*\cap\Delta^*$, અને
$\Gamma^*\cap\Delta^*$ પૂર્ણ છે.

$\Gamma^*$ પૂર્ણ સુસંગત હોવાથી તેને એવો નિદર્શ $\Struct{M}'_1$ મળે
છે જેનું !!{domain}~$\Domain{M'_1}$, !!{constant}sનાં અર્થઘટનો
$\Assign{c}{M'_1}$ એવા અને માત્ર એવા ઘટકોનું બનેલું છે---પૂર્ણતા
પ્રમેયની સાબિતીની જેમ (\olref[fol][com][cth]{thm:completeness}). એ જ
રીતે, $\Delta^*$ને એવો નિદર્શ~$\Struct{M}'_2$ મળે છે જેનું
!!{domain}~$\Domain{M'_2}$, !!{constant}sનાં અર્થઘટનો
$\Assign{c}{M'_2}$થી બનેલું છે.

$\Struct{M'_1}$માંથી $\Lang{L'_1}\setminus\Lang{L'_0}$માંના
!!{constant}s, !!{function}s અને !!{predicate}sનાં અર્થઘટનો કાઢીને
$\Struct{M_1}$ મેળવો, અને $\Struct{M_2}$ પણ એ જ રીતે મેળવો. ત્યારે
$h(\Assign{c}{M'_1})=\Assign{c}{M'_2}$થી વ્યાખ્યાયિત કરેલો નકશો
$h\colon M_1\to M_2$, $\Lang{L}'_0$માં એકરૂપતા છે, કારણ કે બતાવ્યા
પ્રમાણે $\Gamma^*\cap\Delta^*$, $\Lang{L}'_0$માં પૂર્ણ સુસંગત છે.
આનું કારણ એ છે કે કોઈ પણ $\Lang{L}'_0$-!!{sentence} કાં તો $\Gamma^*$
અને~$\Delta^*$ બંનેમાં આવે છે અથવા એકેયમાં આવતું નથી: તેથી
$\Assign{c}{M'_1}\in\Assign{P}{M'_1}$ જો અને માત્ર જો
$\Atom{P}{c}\in\Gamma^*$, જો અને માત્ર જો $\Atom{P}{c}\in\Delta^*$,
જો અને માત્ર જો $\Assign{c}{M'_2}\in\Assign{P}{M'_2}$. એકરૂપતાની
બીજી શરતો પણ એ જ રીતે સ્થાપિત કરી શકાય છે.

હવે !!{language}~$\Lang{L_1}\cup\Lang{L_2}$ માટેનો નિદર્શ
$\Struct{M}$ આ પ્રમાણે વ્યાખ્યાયિત કરીએ:
\begin{enumerate}
\item !!{domain}~$\Domain{M}$ એ $\Domain{M_2}$ જ છે, એટલે બધા
  ઘટકો~$\Assign{c}{M'_2}$નો ગણ;
\item જો કોઈ !!a{predicate}~$P$, $\Lang{L_2}\setminus\Lang{L_1}$માં
  હોય, તો $\Assign{P}{M}=\Assign{P}{M'_2}$;
\item જો કોઈ વિધેય-સંકેત~$P$, $\Lang{L}_1\setminus\Lang{L}_2$માં હોય, તો
  $\Assign{P}{M}=h(\Assign{P}{M'_1})$, એટલે
  $\tuple{\Assign{c_1}{M'_2},\dots,\Assign{c_n}{M'_2}}\in
  \Assign{P}{M}$ જો અને માત્ર જો
  $\tuple{\Assign{c_1}{M'_1},\dots,\Assign{c_n}{M'_1}}\in
  \Assign{P}{M'_1}$.
\item જો કોઈ !!a{predicate}~$P$, $\Lang{L}_0$માં હોય, તો
  $\Assign{P}{M}=\Assign{P}{M'_2}=h(\Assign{P}{M'_1})$.
\item !!^{function}sને, જેમાં !!{constant}s પણ સામેલ છે,
  $\Lang{L}_1\cup\Lang{L}_2$ પર એ જ રીતે સંભાળો.
\end{enumerate}
\sourcecorrection{OLINT-006}{સ્થિર અંગ્રેજી મૂળની
  \texttt{interpolation-proof.tex}, પંક્તિ~173માં
  $P\in\Lang{L}_1\setminus\Lang{L}_2$ હોવા છતાં તેના અર્થઘટનને
  $\Struct{M'_2}$માંથી લઈ $h$ લગાડવામાં આવ્યું છે. $\Struct{M'_2}$
  પાસે $P$નું અર્થઘટન હોવું જરૂરી નથી અને $h$નું ક્ષેત્ર પણ
  $M_1$ છે. તરત પછી આપેલી ``એટલે'' શરત મુજબ અહીં
  $h(\Assign{P}{M'_1})$ પુનઃસ્થાપિત કર્યું છે.}

છેવટે, !!{formula}s પરના આગમનથી બતાવી શકાય છે કે $\Lang{L}_1$નાં
બધાં !!{formula}s માટે $\Struct{M'_1}$ના $\Lang{L}_1$-ન્યૂનરૂપમાં
સંતોષ, $h$ મારફતે $\Struct{M}$ના $\Lang{L}_1$-ન્યૂનરૂપમાં સંતોષને
બરાબર અનુરૂપ છે, અને $\Lang{L}_2$નાં બધાં !!{formula}s માટે
$\Struct{M}$ તથા $\Struct{M'_2}$નાં $\Lang{L}_2$-ન્યૂનરૂપોમાં સંતોષ
સંમત થાય છે. ખાસ કરીને,
$\Struct{M'_1}\Entails\Gamma^*$ તથા $!A\in\Gamma^*$, અને
$\Struct{M'_2}\Entails\Delta^*$ તથા $\lnot !B\in\Delta^*$ હોવાથી
$\Struct{M}\Entails !A$ અને $\Struct{M}\Entails\lnot !B$. તેથી
$\not\Entails !A\lif !B$. આથી ક્રેગના અંતર્વેશન પ્રમેયની સાબિતી
પૂરી થાય છે.
\sourcecorrection{OLINT-007}{સ્થિર અંગ્રેજી મૂળની
  \texttt{interpolation-proof.tex}, પંક્તિઓ~183--187માં
  $\Struct{M}$ને માત્ર $\Lang{L}_1\cup\Lang{L}_2$ માટે વ્યાખ્યાયિત
  કર્યા પછી તે નવા અચળો ધરાવતી વિસ્તૃત ભાષાઓ $\Lang{L}'_1$ અને
  $\Lang{L}'_2$નાં બધાં સૂત્રો પર સંમત થાય છે અને
  $\Gamma^*\cup\Delta^*$ને સંતોષે છે એવો અપરિભાષિત દાવો થાય છે.
  અહીં જરૂરી અને પર્યાપ્ત ન્યૂનરૂપ-દાવો મૂક્યો છે: $h$ પહેલા નિદર્શના
  $\Lang{L}_1$-ભાગને પરિવહિત કરે છે અને બીજાનો
  $\Lang{L}_2$-ભાગ યથાવત્ રહે છે; તેથી મૂળ વાક્યો $!A$ અને
  $\lnot !B$ સાચાં રહે છે.}
\end{proof}

\end{document}
